\documentclass[11pt]{scrartcl}

\usepackage[sexy]{evan}
\usepackage{import}
\usepackage[table]{xcolor}
\usepackage{xfp}
\usepackage{multicol}
\newcommand{\gad}{\textcolor{yellow}{$\bigstar$}}
\newcommand{\bad}{\textcolor{red}{$\bigstar$}}

\definecolor{dcol0}{HTML}{C8E6C9}
\definecolor{dcol1}{HTML}{D4E9B3}
\definecolor{dcol2}{HTML}{E5ED9A}
\definecolor{dcol3}{HTML}{FFF59D}
\definecolor{dcol4}{HTML}{FFE082}
\definecolor{dcol5}{HTML}{FFCC80}
\definecolor{dcol6}{HTML}{FFAB91}
\definecolor{dcol7}{HTML}{F49890}
\definecolor{dcol8}{HTML}{E57373}
\definecolor{dcol9}{HTML}{D32F2F}

\makeatletter
\newcommand{\getcolorname}[1]{dcol#1}
\makeatother

\newcommand{\dif}[1]{%
    \edef\colorindex{\number\fpeval{floor(#1)}}%
    \edef\fulltext{#1}%
    \colorbox{\getcolorname{\colorindex}}{%
        \ifnum\colorindex>8
            \textbf{\textcolor{white}{\,\fulltext\,}}%
        \else
            \textbf{\textcolor{black}{\,\fulltext\,}}%
        \fi
    }%
}
% Variable para dificultad (inicial 0)
\newcommand{\thmdifficulty}{0}

% Comando para asignar dificultad antes del problema
\newcommand{\problemdiff}[1]{\renewcommand{\thmdifficulty}{#1}}

% Estilo del problema que incluye dificultad antes del título
\declaretheoremstyle[
    headfont=\color{blue!40!black}\normalfont\bfseries,
    headformat={%
      \dif{\thmdifficulty}\quad \NAME~\NUMBER\ifx\relax\EMPTY\relax\else\ \NOTE\fi
    },
    postheadspace=1em,
    spaceabove=8pt,
    spacebelow=8pt,
    bodyfont=\normalfont
]{problemstyle}

    \declaretheorem[style=problemstyle,name=Problema,sibling=theorem]{problema}
    \declaretheorem[style=problemstyle,name=Problema,numbered=no]{problema*}

\definecolor{yellow4}{RGB}{255, 206, 0}

\title {Lista Fuentes Diciembres}
\subtitle{Entrenas Nacionales Centros, PAGMOs}
\author{Emmanuel Buenrostro}


\begin{document}
\maketitle
\begin{multicols}{2}
\section{AngleChasing}
\begin{itemize}
    \item \textbf{2.9: } OMM 2011/2
    \item \textbf{2.10: } OMM 2021/4
    \item \textbf{2.11:} OMM 2010/5
    \item \textbf{2.12:} USAMO 2023/1 \gad
    \item \textbf{2.13:} USA TSTST 2023/1 \gad
    \item \textbf{2.14:} EGMO 2015/1
    \item \textbf{2.15:} USAJMO 2018/3
\end{itemize}

\section{Potencia}
\begin{itemize}
    \item \textbf{2.6: } \href{https://medium.com/intuition/a-geometry-problem-solved-using-radical-axes-77b2dbbc282e}{Bekhuz Niyazov}
    \item \textbf{2.9: } USAJMO 2012/1
    \item \textbf{2.11:} USAMO 1998/2
    \item \textbf{2.12:} USAMO 2023/1 \gad 
    \item \textbf{2.13:} Jalisco TST 2025/Final/3, Original mio. 
    \item \textbf{2.14: } Mexico TST 2025 \gad 
    \item \textbf{2.15: } IMO 2009/2
\end{itemize}

\section{Mix}
\begin{itemize}
    \item \textbf{1.3:} ORO 2025/4
    \item \textbf{1.4:} ORO 2024/4
    \item \textbf{1.5:} Jalisco TST 2025/Final/4, Original mio 
    \item \textbf{1.6:} OMCC 2024/4
    \item \textbf{1.7: } OMCC 2020/4
    \item \textbf{1.8: } OMCC 2023/5
    \item \textbf{1.10:} OMCC 2021/2 \gad 
    \item \textbf{1.11: } OMM 2024/4
    \item \textbf{1.12: } OMMock 2025/2
    \item \textbf{1.13:} OMCC 2024/3
    \item \textbf{1.14:} APMO 2015/1
\end{itemize}
\end{multicols}


\end{document}